\documentclass[11pt]{article}
\usepackage{preamble}
\renewcommand{\answer}[1]{}

\begin{document}

\ladderheading{AP Statistics \textperiodcentered\ Topic 2.7 (CED 5.4) \textperiodcentered\ B: 2027-02-11 \textperiodcentered\ E: 2027-04-07}{One Line, Two Model Diagnostics}

\focusbanner{TODAY'S PROBLEM}

Suppose a road-safety lab records braking distance $y$ at speeds $x$ from 10 through 50 mph on one vehicle and dry surface. The line is $\hat y=-10+2.5x$, the correlation is $r=0.997$, and the residual plot is shown. Actual distances at 10, 30, and 50 mph are 19, 61, and 119 feet.\par\smallskip \textbf{Ezra:} The large correlation and residuals on both sides of 0 justify using the line across the tested speeds.\par \textbf{June:} I would use where the residuals occur and the 50-mph error to judge safety, then check any replacement model separately.\par
\begin{center}
\begin{tikzpicture}
\begin{axis}[
  width=0.65\linewidth,
  height=1.75in,
  xmin=5, xmax=55, ymin=-5, ymax=5,
  xlabel={Speed (mph)}, ylabel={Residual (ft)}, grid=major,
  tick label style={font=\small}, label style={font=\small}
]
\addplot[only marks, mark=*, mark size=1.7pt, color=dayblue] coordinates {
  (10,4)
  (15,1)
  (20,-1)
  (25,-2)
  (30,-4)
  (35,-2)
  (40,-1)
  (45,1)
  (50,4)
};
\end{axis}
\end{tikzpicture}
\end{center}
\begin{firsttakebox}{FIRST TAKE --- before the video (5 min)}
Which evidence should control the model judgment? Cite one numerical or graphical feature.
\workspace{0.85in}
\end{firsttakebox}

\begin{callout}{\IconBook\ READING A RESIDUAL PLOT (keep on the page)}{calloutgreen}
A residual is $y-\hat y$. A positive residual means the model underpredicted;
a negative residual means it overpredicted. Plot residuals against $x$ or $\hat y$.
Random scatter around 0 supports a linear form. Curvature, trends, clusters, or
changing spread show structure the line missed. Correlation measures the direction
and strength of linear association, but even a large $|r|$ does not replace checking
residual locations. Check a replacement model with its own residual plot.
\end{callout}

\newpage
\textbf{Finish after the video:}\par\smallskip

\dokbadge{1}~\textbf{(a)} Read the residuals at 10, 30, and 50 mph and state whether the line underpredicts or overpredicts at each speed.\par
\workspace{0.65in}

\dokbadge{2}~\textbf{(b)} Verify those residuals from the equation and actual distances. Describe the full residual plot, including its form and where the line tends to underpredict or overpredict.\par
\workspace{1.35in}

\dokbadge{3}~\textbf{(c)} Evaluate both judgments by reconciling $r=0.997$ with the residual locations. Assess what the large correlation and both residual signs do or do not establish, use the 50-mph error for a safety consequence, specify evidence required before trusting a replacement model, and bound the conclusion to the observed setting.\par
\workspace{1.9in}

\begin{sentenceframebox}\raggedright\textbf{Frame:} The correlation summarizes $\blank[2.1in]$; the residual locations show $\blank[2.4in]$.\end{sentenceframebox}

\begin{sentenceframebox}\raggedright\textbf{Frame:} At 50 mph the line predicts $\blank[0.7in]$ while the actual value is $\blank[0.7in]$, so for safety I would $\blank[2.5in]$.\end{sentenceframebox}

\begin{summaryexitbox}{EXIT --- turn this in}
One thing the video changed about my first take: \hrulefill\par
\workspace{0.4in}
\end{summaryexitbox}

\end{document}
